\section{Pairing obstructions and the four remaining rings}
\label{sec:pairing-frontier-20261027}
All four rings remain unresolved. Every application below uses the established
general-complex existence reduction to a pseudo-unitary Galois representative;
it does not assert unitarity of the original category.

\subsection{Quaternionic invertibles and $R8\_3\_1768$}
For a connected separable algebra in a pseudo-unitary category, its positive
module trace supplies a symmetric separable Frobenius structure. The
bimodule forgetful functor is pivotal Frobenius and preserves a chosen
nondegenerate selfduality pairing with its pivotal sign. Consequently a
negative selfdual bimodule can contain a positive selfdual source simple
only with even multiplicity: its multiplicity space is alternating.
This is pairing transport, not equality of numerical indicators.
See Schaumann, Proposition 5.8 and Theorem 6.6
(\url{https://arxiv.org/abs/1206.5716}), and Mulevi\v cius, Example 4.18
(\url{https://arxiv.org/abs/2206.02611}). The induced positive dimensions
are $\dim U(Q)/\dim A$, which fixes the pivotal normalization.

\begin{theorem}
No category with $GH(C_6,2)$ fusion rules is categorically Morita equivalent
to a categorification of $R8\_3\_1768$, for any associator or algebra.
\end{theorem}
\begin{proof}
Put $u=6+\sqrt{37}$. The positive source second indicators are
$(1,-1,1,1,1,1,1,1)$ and its dimensions are
\[
 (1,1,5+\sqrt{37},6+\sqrt{37},7+\sqrt{37},
                7+\sqrt{37},6+\sqrt{37},5+\sqrt{37}).
\]
Its weighted second-indicator total is $6u$. Such totals are Morita
invariant by the center induction formula. A $GH(C_6,2)$ category has six
selfdual large simples of dimension $u$. Its pointed contribution is
$1+\nu_2(g^3)$. Comparing the irrational coefficient and then the rational
part forces all six large indicators to be $+1$ and $\nu_2(g^3)=-1$.
Thus this Morita dual would contain a negative-indicator invertible.

Choose any simple $m$ in the connecting module and let $A_m$ be its
internal end. For a negative invertible $Q$ in the dual,
$\dim U_m(Q)=\dim A_m$. Pairing parity makes all six nonpointed source
coefficients in $U_m(Q)$ even. Since each has $\sqrt{37}$ coefficient one,
the square-root coefficient in $\dim A_m$ is even for every $m$.

The complete module filter, including every retained incomplete-action
factor, has $21$ possible induction factors. For a root algebra $A$ and
induction matrix $B$, the internal-end dimensions at every module simple
are exactly
\[
 \dim A_m=\frac{(B_m\cdot\boldsymbol d)^2}{\dim A}.
\]
In $17$ factors one of these dimensions has odd square-root coefficient.
They cannot have such a dual. In each of the other four factors a simple
has internal end of dimension one. Connectedness then makes that internal
end the unit, and the internal-end equivalence identifies the module with
the regular module. Its dual has rank eight, not twelve, a contradiction.
The exact filter refactors all $61$ predecessor candidates; the receipt
lists all $17$ odd witnesses and four unit witnesses. It does not discard
the $15$ factors with incomplete actions.
\end{proof}

In particular no algebra $1+g^j\sigma$ in any twisted higher-Haagerup
precursor can model this source. The existence of the auxiliary category
established above remains valid but supplies no model of $R8\_3\_1768$.
The theorem is a construction-family exclusion, not noncategorifiability
of the rank-eight ring itself.

\subsection{Third rotations and $R8\_3\_1154$}
\begin{proposition}
All six noninvertible third indicators equal two in a positive
categorification of $R8\_3\_1154$. Each corresponding trivalent rotation
is identity. All three necessary center patterns still survive.
\end{proposition}
\begin{proof}
The actual rank-ten equivariantization reduction supplies lifts $x,ux$
of $X_4$ and likewise $X_5$, with $u$ central of order two and $x^*=ux$.
For odd $n$ the two character subspaces of the cubic or higher vacuum
space give $\nu_n(X_4)=\nu_n(x)+\nu_n(ux)=2\nu_n(x)$, by spherical
duality. Each lift's cubic vacuum space has dimension one. Thus each of
$\nu_3(X_4),\nu_3(X_5)$ belongs to
$\{2,2\zeta_3,2\zeta_3^2\}$, rather than an arbitrary sum of cube roots.
Imposing these discrete possibilities on all three complete retained
center systems forces both to equal two. Duality and the inner
autoequivalence, followed by a second discrete check, force the other
four noninvertible third indicators to equal two. The $33$ excluded
branches have exact integer separating certificates; all six retained
branch records have exact rational feasibility witnesses. An independent
matrix reconstruction checks every certificate. An order-three rotation
on a two-dimensional space with trace two is identity.
\end{proof}

The integral center is Vec even for an arbitrary complex categorification.
Indeed an integral central simple forgets to $a1+bg$; based centrality and
the noncentrality of $g$ force $b=0$. Full adjoint closure makes a central
object supported on the unit trivial. Hence no Morita representative has
a nontrivial finite-group extension or equivariantization. The auxiliary
rank-ten category has integral center double Semion, whose only nontrivial
Tannakian descent returns to this source.
Direct central condensations from $\mathrm{Fib}^{\boxtimes6}$,
$GH(C_4,1)\boxtimes\mathrm{Fib}^{\boxtimes3}$ and
$GH(C_4,1)^{\boxtimes2}$, each with arbitrary pointed factors, are also
excluded. Dimension and subgroup-index bounds eliminate the first two.
In the last, the required index-three subgroup contains a central $C_4$;
quotienting forces the impossible rank-two rule $x^2=1+4x$.
These are specified central constructions, not all Morita constructions.

\subsection{Simple centralizers for $R8\_3\_928$ and $R8\_3\_941$}
Write $D=144+96\sqrt2+80\sqrt3+56\sqrt6$.
For either source, the only proper nontrivial subcategory is pointed of
rank two. Any central fusion subcategory with a nonintegral object
therefore forgets dominantly to the full source and has dimension at
least $D$. Equality for a braided subcategory would make the noncommutative
source braided, which is impossible.

For R928 its center is $S\boxtimes\mathcal B$, with $S$ Semion and
$\dim\mathcal B=D^2/2$. The only integral simple in $\mathcal B$ is its
unit. A proper nontrivial subcategory and its centralizer would both have
dimension at least $D$, contradicting their product $D^2/2$.
Thus $\mathcal B$ is simple; its possible ranks are $17,20,23$.
The regular-module rigidity theorem and the correspondence between module
categories and Lagrangian algebras imply that $I(1)$ is the unique
Lagrangian algebra. Every connected etale algebra extends to one, by
Witt triviality and lifting local-module algebras. The reverse subcategory
correspondence then gives precisely three etale algebra classes in the
whole center, in a chain $1\subset A\subset I(1)$, with $\dim A=D/2$.
The restriction from $\mathcal B$ is dominant, so its right adjoint gives
that same unique nontrivial $A$ in $\mathcal B$, and
$\mathcal B_A\simeq\mathcal C$, $\mathcal B_A^0\simeq S^{\rm rev}$.
These use DMNO, Lemma 3.5, Proposition 3.16, Corollary 3.30,
Proposition 4.8, Theorem 4.10 and Propositions 5.4--5.8.

For R941 the fermion centralizer $\mathcal E$ has rank $22$, dimension
$D^2/2$, and integral part sVec. A splitting
$\mathcal E=\mathcal M\boxtimes\mathrm{sVec}$ would force the whole
center to have rank $33$ or $44$, instead of $34,37,40,43$ in the
complete retained list. If a proper nonpointed subcategory contains sVec,
it and its centralizer both dominate the source. Their dimension product
is $D^2$, forcing both dimensions to be $D$, the impossible equality
case above. A nonpointed subcategory not containing the fermion would
generate $\mathcal E$ with sVec and be nondegenerate, giving the excluded
splitting. Thus the only subcategories are Vec, sVec and $\mathcal E$.

Finally the conditional $1+X$ Morita construction in the rank-six
manuscript produces precisely R941; it is not forced by equal FP
dimensions. Neither category nor that nonpointed algebra is constructed.
Full proofs, complete finite certificates and audits for this section are
curated in \path{evidence/frontier_20261027/README.md}.
